\documentclass[11pt]{article}
\usepackage{amsmath,amssymb,amsthm}
\usepackage[margin=1.2in]{geometry}
\newtheorem{theorem}{Theorem}
\newtheorem{lemma}[theorem]{Lemma}
\newtheorem{corollary}[theorem]{Corollary}
\newtheorem{remark}[theorem]{Remark}
\title{Erd\H{o}s \#683: reduction to prime-free intervals,\\ and identification with the
Sylvester--Erd\H{o}s problem}
\author{Samuel Lavery}
\date{Draft --- \today}
\begin{document}
\maketitle

\begin{abstract}
\#683 asks whether $P\binom nk\ge\min(n-k+1,k^{1+c})$ for some $c>0$, where $P$ denotes the
largest prime factor.  We prove that the inequality holds automatically --- for every $c>0$
--- whenever $(n-k,n]$ contains a prime, so the problem reduces entirely to prime-free
intervals; by Baker--Harman--Pintz this leaves only $k<n^{0.525+\varepsilon}$.  We further
show that in its main range \#683 \emph{is} the Sylvester--Erd\H{o}s problem on the largest
prime factor of a product of consecutive integers.  \#683 itself is not proved.
\end{abstract}

\section{Setup}
By Kummer, $v_p\binom nk$ is the number of carries in the base-$p$ addition $k+(n-k)$; the
set of primes dividing $\binom nk$ is $\mathcal D(n,k)=\{p:k\not\preceq_p n\}$, where
$k\preceq_p n$ means every base-$p$ digit of $k$ is at most that of $n$.  \#683 asks for
$\max\mathcal D$.  Throughout $1\le k\le n/2$.

\section{The window}

\begin{lemma}[window]\label{lem:window}
Let $p>n/2$ be prime.  Then $v_p\binom nk=1$ if $n-k<p\le n$, and $v_p\binom nk=0$ if
$n/2<p\le n-k$.
\end{lemma}

\begin{proof}
For $p>n/2$ and $m\le n$ we have $p^2>n\ge m$, so $v_p(m!)=\lfloor m/p\rfloor$.  Hence
$v_p\binom nk=\lfloor n/p\rfloor-\lfloor k/p\rfloor-\lfloor(n-k)/p\rfloor$, in which
$\lfloor n/p\rfloor=1$, $\lfloor k/p\rfloor=0$ (as $k\le n/2<p$), and
$\lfloor(n-k)/p\rfloor$ is $1$ or $0$ according as $p\le n-k$ or not.
\end{proof}

\emph{Verified}: no exception, in either direction, for $4\le n<420$ and all $k\le n/2$.

\begin{theorem}[reduction]\label{thm:reduction}
If $(n-k,n]$ contains a prime then $P\binom nk\ge n-k+1$, so the inequality of \#683 holds
for every $c>0$.  Consequently \#683 is equivalent to its restriction to pairs $(n,k)$ for
which $(n-k,n]$ is prime-free.
\end{theorem}

\begin{proof}
A prime $q\in(n-k,n]$ divides $\binom nk$ by Lemma~\ref{lem:window}, so $P\binom nk\ge q>n-k$,
and the minimum is attained by its first branch or exceeded.
\end{proof}

\emph{Measured}: among the $2\,248\,480$ pairs with $10\le n<3000$, $1\le k\le n/2$, the
interval $(n-k,n]$ is prime-free for $13\,909$ --- a proportion of $0.0062$.  So
Theorem~\ref{thm:reduction} settles $99.38\%$ of the range unconditionally.

\begin{corollary}[unconditional range]\label{cor:bhp}
Fix $\varepsilon>0$.  For $n$ large and $n^{0.525+\varepsilon}\le k\le n/2$, the interval
$(n-k,n]$ contains a prime by Baker--Harman--Pintz, so \#683 holds there for every $c>0$.
Only $k<n^{0.525+\varepsilon}$ remains.
\end{corollary}

\emph{Measured}: the longest prime-free $(n-k,n]$ below $n=2\times10^{4}$ has length $51$,
against $n^{0.525}=181$.

\section{Identification with Sylvester--Erd\H{o}s}

\begin{theorem}[refined window]\label{thm:refined}
Let $p>\sqrt n$ be prime with $p>k$.  Then $p\mid\binom nk$ if and only if $p$ has a multiple
in $(n-k,n]$.
\end{theorem}

\begin{proof}
$p^2>n$ gives $v_p(m!)=\lfloor m/p\rfloor$ for $m\le n$, and $k<p$ gives
$\lfloor k/p\rfloor=0$; so $v_p\binom nk=\lfloor n/p\rfloor-\lfloor(n-k)/p\rfloor$, the
number of multiples of $p$ in $(n-k,n]$.
\end{proof}

\emph{Verified}: $7\,327\,302$ triples $(n,k,p)$ with $n<700$ --- no mismatch.

\begin{corollary}\label{cor:sylvester}
If $P^{+}\bigl(\prod_{m=n-k+1}^{n}m\bigr)>\max(\sqrt n,k)$ then
$P\binom nk=P^{+}\bigl(\prod_{m=n-k+1}^{n}m\bigr)$.  Hence in its main range \#683 is
exactly the Sylvester--Erd\H{o}s problem on the largest prime factor of a product of $k$
consecutive integers near $n$.
\end{corollary}

\emph{Verified}: equality with no mismatch for $20\le n<500$, all $k\le n/2$ with the set
nonempty.  \emph{Measured}: over $n<1200$, $3\le k<60$, restricted to pairs where the binding
branch is $k^{1+c}$, the least admissible exponent is $c=0.4650$, at $(n,k)=(10,3)$ with
$P\binom{10}{3}=P(120)=5$; the extremal case is not drifting upward with $n$.

\section{A ledger bound forcing $P\gg k\log k$}

The reduction of \S2 disposes of the range where $(n-k,n]$ contains a prime.  On the
complementary range the size of $\binom nk$ and the ledger bound each rail's contribution,
and the two together bound $P$ from below.

\begin{lemma}[rail ceiling]\label{lem:railceil}
For every prime $p$ and $1\le k\le n$,
\[
  v_p\binom nk\ \le\ \Bigl\lfloor \frac{\log n}{\log p}\Bigr\rfloor+1 .
\]
\end{lemma}

\begin{proof}
By Kummer, $v_p\binom nk$ is the number of carries in the base-$p$ addition $k+(n-k)=n$.
Carries occur only out of positions $0,1,\dots,\lfloor\log_pn\rfloor$, since $n$ has
$\lfloor\log_pn\rfloor+1$ base-$p$ digits and no carry leaves the top position.
\end{proof}

\emph{Verified}: $660\,046$ triples $(n,k,p)$ with $n<260$ --- no violation.

\begin{theorem}[the smoothness inequality]\label{thm:smoothineq}
Let $2\le k\le n/2$ and put $y=P\bigl(\binom nk\bigr)$.  Then
\[
  \frac{k\,\log(n/k)}{\log y}\ \le\ \Omega\!\left(\binom nk\right)\ \le\
  \log n\sum_{p\le y}\frac{1}{\log p} .
\]
\end{theorem}

\begin{proof}
For the left inequality, $\binom nk\ge(n/k)^{k}$, and any integer $N>1$ all of whose prime
factors are at most $y$ satisfies $N\le y^{\Omega(N)}$, whence
$\Omega(N)\ge\log N/\log y$.

For the right inequality, every prime factor of $\binom nk$ is at most $y$, so by
Lemma~\ref{lem:railceil}
\[
  \Omega\!\left(\binom nk\right)=\sum_{p\le y}v_p\binom nk
  \ \le\ \sum_{p\le y}\Bigl(\frac{\log n}{\log p}+1\Bigr),
\]
and the term $\pi(y)$ is absorbed since $\log n/\log p\ge1$ for every $p\le n$.
\end{proof}

\emph{Verified}: $16\,510$ pairs $(n,k)$ with $n<260$ --- both inequalities hold throughout.

\begin{corollary}\label{cor:klogk}
For $k\le\sqrt n$ one has $P\bigl(\binom nk\bigr)\gg k\log k$.
\end{corollary}

\begin{proof}
By the prime number theorem $\sum_{p\le y}1/\log p\sim y/(\log y)^{2}$.  Substituting into
Theorem~\ref{thm:smoothineq} and using $\log(n/k)\ge\tfrac12\log n$ for $k\le\sqrt n$ gives
\[
  \frac{k\log n}{2\log y}\ \le\ (1+o(1))\,\frac{y\log n}{(\log y)^{2}},
  \qquad\text{i.e.}\qquad y\ \ge\ (1+o(1))\,\tfrac12\,k\log y .
\]
Since $y>k$ by Sylvester--Schur, $\log y>\log k$, and the claim follows.
\end{proof}

This recovers Erd\H{o}s's bound $P\binom nk\gg k\log k$ from the ledger alone, with no input
beyond Kummer, Legendre and the prime number theorem.  The gap to \#683 is the passage from
$k\log k$ to $k^{1+c}$: Theorem~\ref{thm:smoothineq} is sharp in the exponent of $y$, so any
improvement must come from a better upper bound on $\Omega\binom nk$ than the rail ceiling
of Lemma~\ref{lem:railceil}, which is attained only when every position carries.

\begin{remark}[register]
\textbf{Proved}: Lemmas and Theorems above, all elementary consequences of Legendre and
Kummer; Corollary~\ref{cor:bhp} additionally cites Baker--Harman--Pintz, and
Corollary~\ref{cor:klogk} the prime number theorem.  Lemma~\ref{lem:railceil},
Theorem~\ref{thm:smoothineq} and Corollary~\ref{cor:klogk} recover Erd\H{o}s's
$P\binom nk\gg k\log k$ from the ledger alone.  None is claimed to
be new --- Lemma~\ref{lem:window} is surely folklore.  What is recorded is the reduction to
prime-free intervals and the identification in Corollary~\ref{cor:sylvester}.
\textbf{Measured}: the $0.62\%$ proportion, $c=0.4650$, the length-$51$ window.
\textbf{Not proved}: \#683.  \textbf{Falsifier}: the least admissible $c$ declining with the
range would indicate no fixed $c$ works; it did not.
\end{remark}
\end{document}
